%% =====================================================================
%%  普及组 CSP-J 2025 初赛模拟卷 2 —— 正文
%%  编译：xelatex csp_mock_paper.tex （需与 cspexam.cls 同目录，跑两遍）
%%  字体依赖：Noto Serif CJK SC / Noto Sans CJK SC / DejaVu Sans Mono
%% =====================================================================
\documentclass{cspexam}

\begin{document}

\examtitle{普及组 CSP-J 2025 初赛模拟卷 2}
\examinfo{120 分钟}{100 分}

% =====================================================================
\bigsection{一、单项选择题}{（共 15 题，每题 2 分，共计 30 分；每题有且仅有一个正确选项）}

\Q 在 C++ 程序中，假设一个字符占用的内存空间是 1 字节，则下列程序中，s 占用的内存空间是（\quad）字节。
\begin{ncode}
char s[] = "hello oiers";
size_t cnt = strlen(s);
cout << cnt << endl;
\end{ncode}
\begin{optrow}
\op A. 10 \op B. 11 \op C. 13 \op D. 12
\end{optrow}

\Q 十六进制数 B2025 转换成八进制数是（\quad）。
\begin{optrow}
\op A. 2620045 \op B. 2004526 \op C. 729125 \op D. 2420045
\end{optrow}

\Q 以下能正确定义二维数组的是（\quad）。
\begin{optgrid}
\op A.\co|int a[3][];|
\op B.\co|int a[][];|
\op C.\co|int a[][4];|
\op D.\co|int a[][2] = {{1,2},{1,2},{3,4}};|
\end{optgrid}

\Q 二进制[10000011]原码和[10000011]补码表示的十进制数值分别是（\quad）。
\begin{optrow}
\op A. $-125$，$-3$ \op B. $-3$，$-125$ \op C. $-3$，$-3$ \op D. $-125$，$-125$
\end{optrow}

\Q 在 C++ 中，下列定义方式中，变量的值不能被修改的是（\quad）。
\begin{optgrid}
\op A.\co|unsigned int a = 5;|
\op B.\co|static double d = 3.14;|
\op C.\co|string s = "ccf csp-j";|
\op D.\co|const char c = 'k';|
\end{optgrid}

\Q 走迷宫的深度优先搜索算法经常用到的数据结构是（\quad）。
\begin{optrow}
\op A. 向量 \op B. 栈 \op C. 链表 \op D. 队列
\end{optrow}

\Q 关于递归，以下叙述中正确的是（\quad）。
\begin{optcol}
\op A. 动态规划算法都是用递归实现的
\op B. 递归比递推更高级，占用的内存空间更少
\op C. A 函数调用 B 函数，B 函数再调用 A 函数不属于递归的一种
\op D. 递归是通过调用自身来求解问题的编程技术
\end{optcol}

\Q 以下不属于计算机输入设备的是（\quad）。
\begin{optrow}
\op A. 扫描仪 \op B. 显示器 \op C. 鼠标 \op D. 麦克风
\end{optrow}

\Q 关于排序算法，下面的说法中正确的是（\quad）。
\begin{optcol}
\op A. 快速排序算法在最坏情况下的时间复杂度是 \(O(n\log n)\)
\op B. 插入排序算法的时间复杂度是 \(O(n\log n)\)
\op C. 归并排序算法的时间复杂度是 \(O(n\log n)\)
\op D. 冒泡排序算法是不稳定的
\end{optcol}

\Q 下列关于 C++ 语言的叙述中不正确的是（\quad）。
\begin{optcol}
\op A. 变量没有定义也能使用
\op B. 变量名不能以数字开头，且中间不能有空格
\op C. 变量名不能和 C++ 语言中的关键字重复
\op D. 变量在定义的时候可以不用赋值
\end{optcol}

\Q 如果 x 和 y 均为 int 类型的变量，下列表达式中能正确判断“x 等于 y”的是（\quad）。
\begin{optgrid}
\op A.\co|(1 == (x / y))|
\op B.\co|(x == (x & y))|
\op C.\co|(0 == (x ^ y))|
\op D.\co!(y == (x | y))!
\end{optgrid}

\Q 在如今的智能互联网时代，AI 如火如荼，除了计算机领域以外，通信领域的技术发展也做出了很大贡献。被称为“通信之父”的是（\quad）。
\begin{optgrid}
\op A. 克劳德·香农
\op B. 莱昂哈德·欧拉
\op C. 约翰·冯·诺依曼
\op D. 戈登·摩尔
\end{optgrid}

\Q 一棵满二叉树的深度为 3（根结点的深度为 1），按照后序遍历的顺序从 1 开始编号，根结点的右子结点的编号是（\quad）。
\begin{optrow}
\op A. 3 \op B. 6 \op C. 7 \op D. 5
\end{optrow}

\Q 三头奶牛 Bessie、Elise 和 Nancy 去参加考试，考场是连续的 6 间牛棚，用栅栏隔开。为了防止作弊，任意两头奶牛都不能在相邻的牛棚，则考场安排共有（\quad）种不同的方法。
\begin{optrow}
\op A. 18 \op B. 24 \op C. 30 \op D. 48
\end{optrow}

\Q 为强化安全意识，某学校准备在某消防月连续 10 天内随机抽取 3 天进行消防紧急疏散演习，抽取的 3 天为连续 3 天的概率为（\quad）。
\begin{optrow}
\op A. 3/10 \op B. 3/20 \op C. 1/15 \op D. 1/18
\end{optrow}

% =====================================================================
\bigsection{二、阅读程序}{（程序输入不超过数组或字符串定义的范围；判断题正确填\checkmark，错误填$\times$；除特殊说明外，判断题每题 1.5 分，选择题每题 3 分，共计 40 分）}

\subprob{（1）}
\begin{ncode}
#include <iostream>
#include <vector>
#include <algorithm>
using namespace std;

using i64 = long long;

int clz(i64 x)
{
    for(int i = 0; i != 64; i++)
        if((x >> (63 - i)) & 1)
            return i;
    return 64;
}

bool cmp(i64 x, i64 y)
{
    if(clz(x) == clz(y))
        return x < y;
    return clz(x) < clz(y);
}

int main()
{
    int n;
    cin >> n;
    vector<int> a(n);
    for(int i = 0; i < n; i++)
        cin >> a[i];
    sort(a.begin(), a.end(), cmp);
    for(int i = 0; i < n; i++)
        cout << a[i] << " \n"[i == n - 1];
    return 0;
}
\end{ncode}

\subsec{判断题}

\Q 若程序输入 5 0 4 2 1 3，则程序输出 4 2 3 1 0。\tfend
\Q 若将第 19 行中的 < 换为 >，则当程序输入 5 0 4 2 1 3 时，程序输出为 4 3 2 1 0。\tfend
\Q 当调用 \co|cmp(3, 3)| 时，函数的返回值为 \co|false|。\tfend

\subsec{选择题}

\Q 若输入 5 4 2 1 3 1，则输出是什么？（\quad）
\begin{optrow}
\op A. 3 4 2 1 1 \op B. 3 2 4 1 1 \op C. 4 3 2 1 1 \op D. 4 2 3 1 1
\end{optrow}

\Q 这个程序实现了什么功能？（\quad）
\begin{optcol}
\op A. 将输入的数组按照二进制位上从左到右第一个 1 前 0 的个数由多到少进行排序
\op B. 将输入的数组按照二进制位上从左到右第一个 1 前 0 的个数由少到多进行排序
\op C. 将输入的数组按照二进制位上从左到右第一个 1 前 0 的个数由多到少进行排序，当 0 的个数相同时，按照原数字由小到大进行排序
\op D. 将输入的数组按照二进制位上从左到右第一个 1 前 0 的个数由少到多进行排序，当 0 的个数相同时，按照原数字由小到大进行排序
\end{optcol}

% ---------------------------------------------------------------------
\subprob{（2）}
\begin{ncode}
#include <iostream>
#include <vector>
#include <algorithm>
#include <set>
#include <string>
using namespace std;

const int inf = 0x3f3f3f3f;

int calc(vector<vector<int>> &grid)
{
    int m = grid.size(), n = grid[0].size();
    vector<vector<int>> dp(m + 1, vector<int>(n + 1, inf));
    dp[0][0] = grid[0][0];
    for(int i = 0; i < m; i++)
        for(int j = 0; j < n; j++)
        {
            if(i > 0)
                dp[i][j]=min(dp[i][j], dp[i-1][j] + grid[i][j]);
            if(j > 0)
                dp[i][j]=min(dp[i][j], dp[i][j-1] + grid[i][j]);
        }
    return dp[m - 1][n - 1];
}

int main()
{
    int m, n;
    cin >> m >> n;
    vector<vector<int>> a(m, vector<int>(n));
    for(int i = 0; i < m; i++)
        for(int j = 0; j < n; j++)
            cin >> a[i][j];
    cout << calc(a) << endl;
    return 0;
}
\end{ncode}

\vspace{0.4em}
\hspace*{\Qindent}假设 \(m\le 100\)，\(n\le 10000\)，完成下面的问题。

\subsec{判断题}

\Q 若输入 2[118;1:3u 3 1 2 3 4 5 6，则输出为 10。\tfend
\Q 计算 dp 数组的时间复杂度为 \(O(n^2)\)。\tfend
\Q （2 分）在 calc 函数中，访问 \co|dp[m][n]| 不会发生越界错误。\tfend

\subsec{选择题}

\Q 当输入的 a 数组为\co|{{1, 3, 1}, {1, 5, 1}, {4, 2, 1}}|时，程序输出为（\quad）。
\begin{optrow}
\op A. 4 \op B. 7 \op C. 6 \op D. 5
\end{optrow}

\Q 若将第 19 行改为 \co|dp[i][j] = min(dp[i][j], dp[i-1][j] - grid[i][j]);|，则当输入的 a 数组为\co|{{1, 2, 3}, {4, 5, 6}}|时，程序的输出为（\quad）。
\begin{optrow}
\op A. $-3$ \op B. $-2$ \op C. $-1$ \op D. 0
\end{optrow}

\Q （4 分）若将第 10 行中的 \& 符号去除，可能出现什么情况？（\quad）
\begin{optgrid}
\op A. dp 数组计算错误
\op B. calc 函数中的 grid 数组和 a 数组不一致
\op C. 无影响
\op D. 发生编译错误
\end{optgrid}

% ---------------------------------------------------------------------
\subprob{（3）}
\begin{ncode}
#include <iostream>
#include <vector>
#include <algorithm>
#include <set>
#include <string>
using namespace std;

const int N = 1010;

vector<int> E[N];
int V[N];
int n;

void add(int x, int y)
{
    E[x].push_back(y);
}

int gcd(int x, int y)
{
    return !y ? x : gcd(y, x % y);
}

void calc(int cur, int fa)
{
    V[cur] = (gcd(cur, fa) != 1);
    for(auto v: E[cur])
    {
        if(v == fa)
            continue;
        calc(v, cur);
        V[cur] += V[v];
    }
    return;
}

int main()
{
    cin >> n;
    for(int i = 1; i < n; i++)
    {
        int x, y;
        cin >> x >> y;
        add(x, y);
        add(y, x);
    }
    calc(1, 1);
    for(int i = 1; i <= n; i++)
        cout << V[i] << " \n"[i == n];
    return 0;
}
\end{ncode}

\vspace{0.4em}
\hspace*{\Qindent}已知 \co|gcd(x, y)| 的时间复杂度为 \(O(\log(\min(x,y)))\)，输入中 \(1\le x,y\le n\) 且 \co|x != y|，回答下面的问题。

\subsec{判断题}

\Q 当输入为 4 1 2 1 3 1 4 时，程序的输出为 0 0 0 0。\tfend
\Q gcd 函数用来计算两个数 x 和 y 的最大公约数。\tfend
\Q 对于树上的一条边\co|(x, y)|，若 x 为 y 的父结点，则必然有 \co|V[x] <= V[y]|。\tfend

\subsec{选择题}

\Q 当输入为 4 1 2 1 3 2 4 时，程序的输出为（\quad）。
\begin{optrow}
\op A. 0 0 0 0 \op B. 1 1 0 1 \op C. 0 0 1 0 \op D. 1 1 1 1
\end{optrow}

\Q （4 分）calc 函数的时间复杂度为（\quad）。
\begin{optrow}
\op A. \(O(n\log n)\) \op B. \(O(n^2)\) \op C. \(O(n)\) \op D. \(O(\log n)\)
\end{optrow}

\Q 若将第 32 行中的代码改为 \co|V[cur] *= V[v]|，则当输入为 4 1 2 1 3 2 4 时，得到的输出为（\quad）。
\begin{optrow}
\op A. 1 1 1 0 \op B. 1 1 0 1 \op C. 0 0 1 0 \op D. 0 0 0 1
\end{optrow}

% =====================================================================
\bigsection{三、完善程序}{（单选题，每小题 3 分，共计 30 分）}

\probdesc{（1）题目描述：}
\hspace*{\Qindent}给定一个长为 n（\(1\le n\le 2\times10^5\)）的数组 a（\(-10^9\le a[i]\le 10^9\)），执行如下操作，直到 a 中只剩下 1 个数：

\hspace*{\Qindent}删除 a[i]。如果 a[i] 左右两边都有数字，则把 a[i$-$1] 和 a[i$+$1] 合成一个数。输出最后剩下的那个数的最大值。
\begin{ncode}
#include <bits/stdc++.h>
using namespace std;

using i64 = long long;

void solve()
{
    int n, flag = 0, mx = (@\Circ{1}@);
    cin >> n;
    vector<int> a(n + 1);
    for(int i = 1; i <= n; i++)
    {
        cin >> a[i];
        if(a[i] < 0)
            flag++;
        mx = max(mx, a[i]);
    }
    if((@\Circ{2}@))
    {
        cout << mx << endl;
        return;
    }
    (@\Circ{3}@) sum1 = 0, sum2 = 0;
    for(int i = 1; i <= n; i += 2)
        sum1 += max(a[i], 0);
    for(int i = 2; i <= n; i += 2)
        (@\Circ{4}@);
    cout << (@\Circ{5}@) << endl;
    return;
}

int main()
{
    int t = 1;
    cin >> t;
    while(t--)
        solve();
}
\end{ncode}

\Q ①处应填（\quad）。
\begin{optrow}
\op A. 0 \op B. 1E9 \op C. $-$1E8 \op D. $-$2E9
\end{optrow}

\Q ②处应填（\quad）。
\begin{optrow}
\op A.\co|flag==n| \op B.\co|flag==0| \op C.\co|flag!=0| \op D.\co|flag!=n|
\end{optrow}

\Q ③处应填（\quad）。
\begin{optrow}
\op A.\co|int| \op B.\co|i64| \op C.\co|i32| \op D.\co|unsigned int|
\end{optrow}

\Q ④处应填（\quad）。
\begin{optgrid}
\op A.\co|sum2 += max(a[i], 0)|
\op B.\co|sum2 += min(a[i], 0)|
\op C.\co|sum2 -= min(a[i], 0)|
\op D.\co|sum2 -= max(a[i], 0)|
\end{optgrid}

\Q ⑤处应填（\quad）。
\begin{optgrid}
\op A.\co|min(sum1, sum2)|
\op B.\co|max(sum1, sum2)|
\op C.\co|sum1 + sum2|
\op D.\co|sum1 - sum2|
\end{optgrid}

% ---------------------------------------------------------------------
\probdesc{（2）题目描述：}
\hspace*{\Qindent}给定一个字符串 t 和一个字符串列表 s 作为字典。保证 s 中的字符串互不相同，且 t 和 s[i] 中均只包含小写英文字母。

\hspace*{\Qindent}如果可以利用字典中出现的一个或多个单词拼接出 t，则返回 true。注意：不要求字典中出现的单词全部使用，并且字典中的单词可以重复使用。

\hspace*{\Qindent}数据限制：\(1\le\) t.length()\,\(\le 300\)，\(1\le\) s.size()\,\(\le 1000\)，\(1\le\) s[i].length()\,\(\le 20\)。
\begin{ncode}
#include <iostream>
#include <vector>
#include <algorithm>
#include <set>
#include <string>
using namespace std;

const int N = 1010;

int n, mx, m;
vector<string> s;
vector<int> mem;
string t;
set<string> st;

int dfs(int i)
{
    if(i == 0)
        return 1;
    if((@\Circ{1}@))
        return mem[i];
    for(int j = i - 1; j >= max(i - mx, 0); j--)
        if(st.find((@\Circ{2}@)) != st.end() && dfs(j))
            return mem[i] = 1;
    return (@\Circ{3}@);

}

int main()
{
    cin >> n;
    s.resize(n);
    for(int i = 0; i < n; i++)
    {
        cin >> s[i];
        mx = max(mx, (@\Circ{4}@));
    }
    st = set<string>(s.begin(), s.end());
    cin >> t;
    m = (int)t.length();
    mem.resize(m + 1, -1);
    if((@\Circ{5}@))
        cout << "Yes\n";
    else
        cout << "No\n";

    return 0;
}
\end{ncode}

\Q ①处应填（\quad）。
\begin{optgrid}
\op A.\co|mem[i] != -1|
\op B.\co|mem[i] == -1|
\op C.\co|mem[i] == 0|
\op D.\co|mem[i] != 0|
\end{optgrid}

\Q ②处应填（\quad）。
\begin{optgrid}
\op A.\co|t.substr(i, j - i)|
\op B.\co|t.substr(j, i - j)|
\op C.\co|t.substr(j)|
\op D.\co|t.substr(i)|
\end{optgrid}

\Q ③处应填（\quad）。
\begin{optgrid}
\op A.\co|1|
\op B.\co|mem[i] = 1|
\op C.\co|0|
\op D.\co|mem[i] = 0|
\end{optgrid}

\Q ④处应填（\quad）。
\begin{optgrid}
\op A.\co|s[i].length()|
\op B.\co|(int)s[i].length()|
\op C.\co|s[i].length() - 1|
\op D.\co|(int)s[i].length() - 1|
\end{optgrid}

\Q ⑤处应填（\quad）。
\begin{optrow}
\op A.\co|!dfs(m)| \op B.\co|!dfs(n)| \op C.\co|dfs(m)| \op D.\co|dfs(n)|
\end{optrow}

\answersheet
\end{document}

